\subsection{5. Clustering huge protein sequence sets in linear time}
\textbf{OSTI:} \texttt{1624105}\quad\textbf{Authors:} Steinegger \& Soeding (2018)\quad\textbf{Domain:} Computational biology / algorithms\\
\scorebadge{Coverage}{8}\quad\scorebadge{Agreement}{9}\quad\textbf{Total:} 17/20\par\smallskip
\noindent\textbf{Coverage rationale.} From-scratch Python reimplementation of Linclust's three-stage algorithm (bottom-m k-mer sketch, bucketed center assignment, ungapped identity verification) with a matched O(N log N) complexity test. Production MMseqs2 binary was not rerun on UniRef-scale datasets (memory footprint, cluster-count and compression ratios on \textgreater{}1e8 sequences) due to the 2-hour compute budget; we benchmarked up to N=64k synthetic proteins.\par
\noindent\textbf{Agreement rationale.} Algorithmic claim reproduces directly: measured log-log slope 1.31 (linclust-py) vs 2.08 (naive O(N\textasciicircum{}2)) --- the paper's linear-time signature. \textasciitilde{}240x speedup at N=4000 and F1 \textgreater{}= 0.987 vs ground truth both match the paper's qualitative performance claims (precision + linear scaling). Absolute wall-time numbers will differ from the paper's C++ MMseqs2 code, but the scaling exponents agree to within expected constant factors.\par\smallskip
\noindent\textbf{What reproduced:}
\begin{itemize}[leftmargin=1.4em,itemsep=0.2em,topsep=0.2em]
\item Bottom-m k-mer sketching as the locality-sensitive primitive
\item Bucketed center-assignment strategy
\item Ungapped identity verification stage
\item Near-linear scaling empirically confirmed on synthetic proteins up to N=64k
\item Cluster-quality F1 \textgreater{}= 0.987 vs ground truth
\item Scaling and quality figures (log-log and linear axes)
\end{itemize}\noindent\textbf{What missing:}
\begin{itemize}[leftmargin=1.4em,itemsep=0.2em,topsep=0.2em]
\item Production MMseqs2 C++ binary run on UniRef50/UniRef100 (1e8-1e9 sequences)
\item Paper's reported wall-time and memory footprint on 1.6B sequences
\item Cluster count / compression-ratio comparison with CD-HIT / UCLUST / kClust
\item Pre-filter vs ungapped-alignment vs Smith-Waterman sensitivity curves
\item Cascaded clustering (linclust -\textgreater{} mmseqs2 iterative) comparison
\end{itemize}\noindent\textbf{Follow-on research questions:}
\begin{enumerate}[leftmargin=1.6em,itemsep=0.2em,topsep=0.2em,label=Q\arabic*.]
\item How does Linclust's clustering quality (F1 vs curated reference families) degrade as intra-family identity is pushed from 50\% down to 30\% on a controlled synthetic benchmark
\item Can the bottom-m sketch be replaced with a learned MinHash via a small Transformer embedding, and does that improve remote-homolog recall at matched runtime
\item What is the GPU speedup achievable by porting the bucketed center-assignment stage to CUDA (hash-table on device memory), and does the memory bandwidth bound dominate
\item How sensitive is the linear-time scaling exponent to the k-mer length k and the sketch width m, and is there a regime where the method silently becomes O(N log N)
\item When inputs are metagenome-scale (\textgreater{}5e9 sequences with heavy sequence redundancy), does Linclust's single-pass center selection still produce stable cluster boundaries under random input shuffling
\end{enumerate}

